\begin{abstract}
We formalize the median dynamics of Doerr, Goldberg, Minder, Sauerwald and Scheideler (SPAA
2011): every node adopts the median of its own value and the values of two random nodes. The
median commutes with monotone maps, so thresholding the process at any value gives the binary
median process, i.e.\ 2-Choices, driven by the same samples. For two values we prove consensus
from a gap of order $\sqrt{n\log n}$ within $O(\log n)$ rounds with high probability, then
consensus from any binary configuration and, by a union bound over the thresholds, from any
configuration (Theorem~1 of the paper, without adversary). With an odd number of equally
supported values, the middle one wins within $O(\log m+\log\log n)$ rounds (the odd case of
Theorem~21).
\end{abstract}

\section{The model}
\begin{definition}[Median dynamics]\label{def:median}
\lean{Median.med3,Median.Config,Median.Round,Median.step,Median.run,Median.ones,Median.Consensus,Median.notConsensus,Median.kernel}\leanok
Values lie in a linear order. In a round every node samples two nodes uniformly at random
(independently, with replacement) and adopts the median of its own value and the two sampled
values. With two values this is 2-Choices: a node keeps its value unless both samples hold the
other one.
\end{definition}

\section{Structure}
\begin{theorem}[Threshold reduction]\label{thm:threshold}
\lean{Median.med3_monotone,Median.threshold_step,Median.threshold_run}\leanok\uses{def:median}
The median commutes with monotone maps; hence for every threshold $\theta$ the indicator of
$\{x_v\ge\theta\}$ evolves, round by round and for the same samples, as the binary median process.
\end{theorem}
\begin{proof}\leanok
$\mathrm{med}(a,b,c)=\max(\min(a,b),\min(\max(a,b),c))$, and monotone maps commute with $\min$ and $\max$.
\end{proof}
\begin{theorem}[Validity, range, consensus]\label{thm:basic}
\lean{Median.step_mem,Median.step_mem_Icc,Median.step_of_consensus,Median.absorbed}\leanok\uses{def:median}
Every new value is a current value, values stay within any interval containing all current
values, consensus configurations are fixed, and consensus is reached almost surely.
\end{theorem}
\begin{proof}\leanok
The median is one of its arguments; the round in which every node samples node $0$ twice gives
consensus and has positive probability; conclude with the shared absorption theorem.
\end{proof}
\begin{theorem}[Binary expectation]\label{thm:expect}
\lean{Median.expected_ones}\leanok\uses{def:median}
With two values and a fraction $p$ of \texttt{true}, the expected number of \texttt{true} after one
round is $n(3p^2-2p^3)$, the same cubic as for 3-Majority.
\end{theorem}
\begin{proof}\leanok
A \texttt{true} node stays unless both samples are \texttt{false}, a \texttt{false} node switches iff
both samples are \texttt{true}.
\end{proof}

\section{Consensus from a vanishing bias}
\begin{theorem}[2-Choices, w.h.p.]\label{thm:whp}
\lean{Median.binary_consensus,Median.consensus_whp}\leanok\uses{thm:expect}
If $\log n\ge 128$ and the \texttt{true} nodes outnumber the \texttt{false} ones by at least
$128\sqrt{n\log n}$, all nodes hold \texttt{true} after $\lceil 128\log n\rceil$ rounds with
probability at least $1-128/n$.
\end{theorem}
\begin{proof}\leanok
Phases as in the plurality package: while it is small, the gap grows by a factor $5/4$ per phase
except with probability $O(n^{-2})$ (Bernstein); then the minority shrinks by $7/8$ per phase down
to $O(\log n)$; a final round reaches consensus except with probability $O(n^{-1/2})$ (Markov). The
phases are combined with the nested-phases lemma of the shared layer. The threshold
$\log n\ge128$, i.e.\ $n\ge e^{128}$, comes from these crude constants.
\end{proof}

\section{Consensus from any configuration}
\begin{theorem}[2-Choices from any start]\label{thm:binary-any}
\lean{Median.binary_any_start}\leanok\uses{thm:whp}
There is a constant $C>0$ such that, whenever $\log n\ge C$, from every binary configuration
(possibly perfectly balanced) all nodes agree after $\lceil C\log n\rceil$ rounds except with
probability at most $C/n$.
\end{theorem}
\begin{proof}\leanok
Take $C=2^{18}$ and write $g$ for the number of \texttt{true} nodes minus the number of
\texttt{false} nodes. \emph{Escape:} the potential $e^{-|g|/(256\sqrt n)}$ contracts in
expectation by $e^{-1/65536}$ per round, up to an additive error $e^{1/131072-\sqrt n/512}$:
away from balance by Hoeffding's lemma applied at every node, near balance because the next gap
has variance at least $n/10$, so by Paley--Zygmund it leaves $[-\sqrt n/4,\sqrt n/4]$ with
probability at least $9/64$. Iterating over $\lceil 2^{17}\log n\rceil$ rounds and applying
Markov's inequality, $|g|\ge 128\sqrt{n\log n}$ except with probability $2/n$.
\emph{Consensus:} Theorem~\ref{thm:whp}, applied to the configuration or to its flip, finishes
within $\lceil 128\log n\rceil$ rounds except with probability $128/n$.
\end{proof}
\begin{theorem}[Reduction to two values]\label{thm:of-binary}
\lean{Median.consensus_of_binary}\leanok\uses{thm:threshold}
If every binary configuration fails to reach consensus within $T$ rounds with probability at
most $\varepsilon$, then a configuration with $m$ distinct values fails with probability at most
$(m-1)\varepsilon$.
\end{theorem}
\begin{proof}\leanok
If a run has not reached consensus, two nodes hold values $a<b$, and the threshold configuration
$\{x_v\ge b\}$, which follows the same rounds by Theorem~\ref{thm:threshold}, has not reached
consensus either. Take a union bound over the $m-1$ thresholds above the minimum value.
\end{proof}
\begin{theorem}[Consensus from any configuration]\label{thm:any}
\lean{Median.median_consensus_any}\leanok\uses{thm:binary-any,thm:of-binary}
There is a constant $C>0$ such that, whenever $\log n\ge C$, from every configuration with values
in a linear order (any number of distinct values) all nodes agree after $\lceil C\log n\rceil$
rounds except with probability at most $C/n$. This is Theorem~1 of Doerr et al.\ without the
adversary.
\end{theorem}
\begin{proof}\leanok
Consensus is absorbing, so three blocks of $\lceil C_0\log n\rceil$ rounds turn the binary failure
bound $C_0/n$ of Theorem~\ref{thm:binary-any} into $(C_0/n)^3$. The union bound of
Theorem~\ref{thm:of-binary} over at most $n-1$ thresholds gives $n(C_0/n)^3\le (3C_0+3)/n$, since
$C_0^3\le n$ when $\log n\ge 3C_0+3$.
\end{proof}

\section{Many values: the middle one wins fast}
\begin{theorem}[Fast binary consensus from a large gap]\label{thm:binary-fast}
\lean{Median.binary_consensus_fast}\leanok\uses{def:median}
There is a constant $C>0$ such that, whenever $\log n\ge C$, if the \texttt{true} nodes outnumber
the \texttt{false} ones by at least $\Delta\ge C\sqrt{n\log n}$, all nodes hold \texttt{true} after
$\lceil C(\log(n/\Delta)+\log\log n)\rceil$ rounds with probability at least $1-C/n$.
\end{theorem}
\begin{proof}\leanok
Take $C=128$ and chain one-round moves, each failing with probability at most $n^{-2}$, through
four segments: the gap grows by a factor $5/4$ per round until the minority is below $n/4$
($O(\log(n/\Delta))$ rounds); six rounds bring the minority below $n/8$; it then shrinks
quadratically, $m\mapsto\approx 3m^2/n$ (Bernstein), down to $512\log n$ ($O(\log\log n)$
rounds); eight more rounds give consensus. Consensus is absorbing, so the remaining rounds are
harmless, and a union bound over the rounds gives the failure probability.
\end{proof}
\begin{theorem}[Margins around a value]\label{thm:median-fast}
\lean{Median.median_consensus_fast}\leanok\uses{thm:threshold,thm:binary-fast}
There is a constant $C>0$ such that, whenever $\log n\ge C$: if, for a value $v$, the nodes holding a value
$\ge v$ outnumber those holding a value $<v$ by at least $\Delta\ge C\sqrt{n\log n}$, and
symmetrically the nodes holding a value $\le v$ outnumber those holding a value $>v$ by at least
$\Delta$, then all nodes hold $v$ after $\lceil C(\log(n/\Delta)+\log\log n)\rceil$ rounds with
probability at least $1-2C/n$.
\end{theorem}
\begin{proof}\leanok
The indicators of $\{x_u\ge v\}$ and of $\{x_u\le v\}$ both evolve as 2-Choices with the same
samples (the second because the median is self-dual), with gaps the two margins. When both reach
consensus on \texttt{true}, all nodes hold $v$; conclude with Theorem~\ref{thm:binary-fast}
(with the same $C=128$) and a union bound.
\end{proof}
\begin{theorem}[An odd number of equally supported values]\label{thm:odd-split}
\lean{Median.odd_split_consensus}\leanok\uses{thm:median-fast}
There is a constant $C>0$ such that, whenever $\log n\ge C$: if the nodes hold $2k+1$ distinct
values, each held by $n/(2k+1)$ nodes, and $C(2k+1)\sqrt{n\log n}\le n$, then the middle value is
held by all nodes after $\lceil C(\log(2k+1)+\log\log n)\rceil$ rounds with probability at least
$1-C/n$. This is the deterministic core of the odd case of Theorem~21 of Doerr et al.
\end{theorem}
\begin{proof}\leanok
Both margins around the middle value equal $n/(2k+1)\ge C_0\sqrt{n\log n}$, and
$\log(n/\Delta)=\log(2k+1)$; apply Theorem~\ref{thm:median-fast} with $C=2C_0$.
\end{proof}
